% La proyección operacional tiene recursos separados del almacén de autoridad.
\begin{figure}[htbp]
  \centering
  \begin{tikzpicture}[
    bloque/.style={draw=SynNavyLight,fill=SynSurface,rounded corners=1mm,
      align=left,font=\fontsize{9}{11}\selectfont,text width=69mm,inner sep=2mm},
    banda/.style={bloque,text width=148mm,fill=SynNavy,text=SynWhite},
    enlace/.style={-{Latex},draw=SynNavyLight,thick}]
    \node[banda] (entrada) {\textbf{AWS São Paulo, sa-east-1}\\Route 53,
      CloudFront, WAF/Shield, HTTPS y API Gateway REST Regional; Cognito OIDC/SAML y MFA.
      VPC Link V2 hacia ALB interno; aplicación y datos privados.};
    \node[bloque,below=7mm of entrada.south west,anchor=north west] (aza)
      {\textbf{Zona de disponibilidad A}\\ECS/Fargate: núcleos Operacional,
       de Servicios y Administrativo; procesos especializados.};
    \node[bloque,below=7mm of entrada.south east,anchor=north east] (azb)
      {\textbf{Zona de disponibilidad B}\\Tareas redundantes de los mismos
       núcleos; escalado y reemplazo por comprobaciones de salud.};
    \node[bloque,below=7mm of aza] (autoridad)
      {\textbf{RDS PostgreSQL 18 Multi-AZ}\\Primario y standby entre zonas;
       autoridad transaccional. 200 GB iniciales, ampliables a 1 TB.};
    \node[bloque,below=7mm of azb] (lectura)
      {\textbf{PostgreSQL de proyecciones}\\Multi-AZ; recursos separados y eventos del
       propietario. Reserva lógica 15 GB; prueba a 45 GB. Vista reconstruible,
       sin autoridad de confirmación.};
    \node[bloque,below=7mm of autoridad] (eventos)
      {\textbf{Eventos e integración}\\EventBridge, SQS/SNS; adaptadores
       versionados. Notificaciones, reintentos y cola de fallos.
       Webhooks por API y salida NAT redundante.};
    \node[bloque,below=7mm of lectura] (objetos)
      {\textbf{Objetos y explotación histórica}\\S3 con versiones y Object
       Lock; Parquet/Athena desacoplado. Textract propone campos con
       revisión humana y alternativa manual.};
    \node[bloque,below=7mm of eventos] (recinto)
      {\textbf{Recinto}\\Dos túneles IPsec y HTTPS/mTLS; contratos, hechos,
       secuencias y acuses. Célula Local con autoridad operacional
       y objetivo de 24 h sin WAN; su cobertura habilitante se desarrolla
       en el modelo de acceso local.};
    \node[bloque,below=7mm of objetos] (dr)
      {\textbf{AWS México Central}\\Réplica PostgreSQL asíncrona, S3 CRR,
       ECR e infraestructura reproducible. Promoción autorizada de
       una única autoridad; warm standby.};
    \node[banda,below=7mm of recinto.south west,anchor=north west] (transversal)
      {IAM, KMS y Secrets Manager; OpenTelemetry/CloudWatch;
       CloudTrail, GuardDuty Runtime Monitoring, Config, Security Hub e Inspector.
       Registros protegidos en cuenta separada; OpenTofu y artefactos ECR.};
    \draw[enlace] (entrada.south -| aza.north) -- (aza.north);
    \draw[enlace] (entrada.south -| azb.north) -- (azb.north);
    \draw[enlace] (aza) -- (autoridad);
    \draw[enlace] (azb) -- (lectura);
    \draw[enlace] (autoridad) -- (eventos);
    \draw[enlace] (eventos.east) -- node[above,font=\fontsize{9}{11}\selectfont]
       {contratos} (objetos.west);
    \draw[enlace] (eventos.north east) -- (lectura.south west);
    \draw[enlace] (eventos) -- (recinto);
    \draw[enlace] (objetos) -- (dr);
  \end{tikzpicture}
  \caption{Despliegue físico detallado de la Zona 1: AWS São Paulo}
  \label{fig:4-9}
  \FuenteFigura{Elaboración de Synaptix; SD4, sección 4.2, y SD5, secciones 5.2.3 y 5.4. Las reservas de capacidad son hipótesis de diseño.}
\end{figure}
